\def\stat{adam-volkov}





\def\tit{ПОДХОД К~СВЯЗЫВАНИЮ ЗАПИСЕЙ В~ТЕХНОЛОГИИ~ПОДДЕРЖКИ\\


КОНКРЕТНО-ИСТОРИЧЕСКИХ ИССЛЕДОВАНИЙ,\\ ОСНОВАННЫЙ НА~НЕЧЕТКИХ 


МНОЖЕСТВАХ}





\def\titkol{Подход к~связыванию записей в~технологии ПКИИ,


%поддержки  конкретно-исторических исследований, 


основанный на~нечетких  множествах}





\def\aut{И.\,М.~Адамович$^1$, О.\,И.~Волков$^2$}





\def\autkol{И.\,М.~Адамович, О.\,И.~Волков}





\titel{\tit}{\aut}{\autkol}{\titkol}





\index{И.\,М.~Адамович$^1$, О.\,И.~Волков$^2$}


\index{I.\,M.~Adamovich and O.\,I.~Volkov}





%{\renewcommand{\thefootnote}{\fnsymbol{footnote}}


%\footnotetext[1]


%{Статья публикуется по представлению программного комитета VI Международной научно-


%образовательных и~ведомственных информационных сетей>>.}} 





\renewcommand{\thefootnote}{\arabic{footnote}}


\footnotetext[1]{Федеральный исследовательский центр <<Информатика и~управление>> Российской академии 


наук, \mbox{Adam@amsd.com}}


\footnotetext[2]{Федеральный исследовательский центр <<Информатика и~управление>> Российской академии 


наук, \mbox{Volkov@amsd.com}}





 \label{st\stat}





 \thispagestyle{headings}


 


 \vspace*{-10pt} 


 





  


  \Abst{Статья посвящена дальнейшему развитию распределенной технологии 


поддержки конкретно-исторических исследований (ПКИИ), основанной на принципах 


краудсорсинга и~ориентированной на широкий круг не относящихся к~профессиональным 


историкам и~биографам пользователей. Развитие осуществляется за счет модификации 


алгоритма автоматического связывания записей при обработке номинативных документов 


семейной структуры, опирающегося на метод восстановления истории семей, успешно 


применяемый в~исторической демографии, и~структуру графов родственных связей 


в~семье. Модификация заключается в~опоре на теорию нечетких множеств 


и~использовании мер сходства для принятия решения о~связывании записей в~случае, когда 


алгоритмы метода восстановления истории семей не позволяют сделать однозначный 


вывод о~связанности. Описывается методика построения функций принадлежности 


и~обосновывается использование меры сходства, основанной на мерах сходства по 


Лукасевичу и~по Заде.}


  


  \KW{конкретно-историческое исследование; распределенная технология; связывание 


записей; источники семейной структуры; нечеткие множества}





\DOI{10.14357\!/\!08696527220213} 





\vskip 10pt plus 9pt minus 6pt





 \thispagestyle{headings}


 


 \vspace*{-8pt}


  


\section{Введение}





\vspace*{-2pt}





  Поддержка конкретно-исторических исследований стала одной из 


актуальных задач современности, что обусловлено вовлечением 


в~исследовательский процесс не только членов профессионального 


исторического сообщества, но и~самых широких слоев непрофессионалов 


в~связи со все возрастающим интересом к~частной, семейной истории~[1].


  


  В~[2, 3] описана разработанная в~ФИЦ ИУ РАН распределенная 


технология ПКИИ, основанная на принципах краудсорсинга (мобилизации 


ресурсов широкого круга добровольцев посредством информационных 


технологий), и~создан ряд программных средств ее поддержки. 


  


  Вся совокупность исторических источников, содержащих основной объем 


исходных данных для кон\-крет\-но-ис\-то\-ри\-че\-ско\-го исследования, 


делится на три основные группы: номинативные (содержащие персональную 


информацию с~указанием имени собственного), статистические 


и~нарративные (повествовательные) источники~[4]. Наибольший интерес 


для исследователя представляют номинативные источники, 


подразделяющиеся на источники персональной структуры и~источники 


семейной структуры~[5].


  


  При обработке номинативных источников одной из важнейших 


и~сложнейших задач является поиск соответствий и~взаимоувязывание 


данных. Связывание записей (Record Linkage)~--- классическая проблема 


в~исторической информатике. Несмотря на наличие значительных 


негативных особенностей номинативных исторических источников, задача 


автоматизированного связывания записей для информации, полученной из 


таких источников, довольно успешно решается в~рамках метода 


восстановления истории семей (ВИС), используемого в~исторической 


демографии~[4].


  


  В задачах связывания записей используются детерминированные 


и~вероятностные алгоритмы. Детерминированные алгоритмы основаны на 


полном совпадении атрибутов записей, и~для целей ВИС они применяться 


практически не могут. Более предпочтительной альтернативой являются 


вероятностные алгоритмы. В~этом случае для атрибутов задаются весовые 


коэффициенты, определяющие влияние на итоговую оценку вероятности 


того, что оцениваемые записи относятся к~одному и~тому же объекту. Пары 


записей с~вероятностями выше некоторого порога считаются связанными, 


а~пары с~вероятностями ниже другого порога считаются несвязанными. 


Пары, которые попадают между этими двумя пороговыми значениями, 


считаются кандидатами на связывание и~обрабатываются отдельно 


(например, вручную).


  


  Опыт реальных исследований показывает, что наибольших успехов 


в~об\-ласти автоматизации связывания записей исторического регистра 


в~рамках ВИС удается добиться при обработке метрических записей 


и~подобных персональных источников демографических событий, но при 


обработке источников семейной структуры возникают проблемы~[6]. Из 


этого следует, что имеется потребность в~механизмах автоматического 


связывания для семейных структур в~ситуации, когда стандартные 


алгоритмы ВИС не смогли установить связь ни для одного члена семьи.





\section{Структурный подход к~связыванию семей}





  Связывание записей в~технологии ПКИИ осуществляется сравнением 


персон, информация о которых уже зафиксирована в~семантической сети 


технологии, с~персонами, упоминаемыми в~некотором номинативном 


историческом документе. В~[7] описан подход к~связыванию записей 


в~ситуации, когда при обработке номинативных источников семейной 


структуры автоматическими алгоритмами связывания метод ВИС не дает 


уверенного ответа о связанности, т.\,е.\ нет пар с~вероятностью связанности, 


достаточной, чтобы считать их определенно связанными, но имеются пары 


с~вероятностью связанности выше порога определенной несвязанности. Для 


этих пар будем говорить о наличии потенциальной связи (P-связи). 


  


  Поскольку родственные связи образуют направленный граф 


с~именованными связями и~вершинами, соответствующими персонам, 


алгоритм опирается на оценку сходства структур таких графов. Для 


сравнения структур алгоритм учитывает:


  \begin{enumerate}[(1)]


  \item  объем потенциальных совпадений в~структурах графов (доля узлов 


с~P-свя\-зями);


  \item  объем признанных на основании экспертной оценки естественными 


(легкообъяснимыми) несовпадений в~структурах графов (доля 


соответствующих узлов);


  \item объем труднообъяснимых несовпадений в~структурах графов (доля 


прочих узлов).


  \end{enumerate}


  


  Формула вычисления вероятности признания графов сходными 


представляет собой взвешенную разность доли потенциальных совпадений 


и~доли труднообъяснимых несовпадений. Весовые коэффициенты 


определяются по результатам испытаний алгоритма в~реальных  


кон\-крет\-но-ис\-то\-ри\-че\-ских исследованиях. Рост объема совпадений 


увеличивает вероятность признания графов сходными, а~рост объема 


труднообъяснимых несовпадений~--- уменьшает. При превышении 


с~вычисленной по соответствующей формуле вероятностью некоторого 


порогового значения записи с~P-свя\-зя\-ми признаются связанными. 


  


  Несмотря на то что данный подход существенно дополняет и~развивает 


технологию ПКИИ, он не свободен от ряда недостатков. Основные из них~--- 


неудовлетворительная точ\-ность некоторых критериев и~субъективность 


в~определении ряда па\-ра\-мет\-ров алгоритма.


  


  Недостаточная точность вытекает из дихотомического подхода 


к~определению естест\-вен\-ности\!/\!труд\-но\-объ\-яс\-ни\-мости 


несовпадений в~структурах графов. В~реальности мнения экспертов часто 


оцениваются по порядковой шкале~[8].


  


  Любое автоматическое принятие решения невозможно без задания 


некоторых пороговых значений исследователем исходя из его субъективных 


предпочтений и~опыта. Но рассматриваемый алгоритм имеет неоправданно 


большое число таких задаваемых пользователем параметров:


  \begin{itemize}


\item параметры алгоритмов ВИС, позволяющие считать записи 


связанными, потенциально связанными или несвязанными;


\item весовые коэффициенты формулы вычисления вероятности 


признания графов сходными;


\item пороговое значение вычисленной вероятности признания графов 


сходными, превышение которого позволяет признать записи  


с~P-свя\-зя\-ми связанными.


\end{itemize}





  Из недостаточной точности и~субъективности рассмотренного метода 


вытекает необходимость в~построении нового алгоритма связывания записей 


в~технологии ПКИИ при обработке номинативных источников семейной 


структуры, опи\-ра\-юще\-го\-ся на алгоритмы ВИС и~свободного от этих 


недостатков.


  


\section{Подход на основе нечетких множеств}





  Предлагается новый подход к~решению задачи связывания записей при 


обработке документов семейной структуры в~случае, когда метод ВИС не 


дает однозначного ответа. Подход основан на нечетких множествах.


  


  Под нечетким множеством~$A$ понимается совокупность упорядоченных 


пар, составленных из элементов~$x$ универсального множества~$X$ 


и~со\-от\-вет\-ст\-ву\-ющих степеней принадлежности $\mu_A(x)$:  


$A\hm= \left\{ (x,\mu_A(x))\vert x\hm\in X\right\}$, где $\mu_A(x)$~--- 


функция принадлежности, принимающая, как правило, значения в~диапазоне 


$[0,1]$ и~указывающая, в~какой степени (мере) элемент~$x$ принадлежит 


нечеткому множеству~$A$~[9].


  


  Для нечетких множеств определено отношение сходства. Среди большого 


числа мер сходства наиболее широкое распространение получили:


  \begin{enumerate}[(1)]


  \item  мера сходства по Лукасевичу:


  


  \noindent


  $$


  m_L\left( \tilde{A},\tilde{B}\right) =1-\fr{\sum\nolimits^n_{i=1} \vert 


\mu_{\tilde{A}}(x_i)-\mu_{\tilde{B}}(x_i)\vert}{n}\,;


  $$


  \item мера сходства по Заде:


  


\vspace*{-3pt}





  \noindent


\begin{multline*}


  m_Z\left(\tilde{A},\tilde{B}\right) ={}\\


  {}=\fr{\sum\nolimits^n_{i=1} \min \left(\max 


(1-\mu_{\tilde{A}}(x_i), \mu_{\tilde{B}}(x_i)), \max (\mu_{\tilde{A}}(x_i), 1-


\mu_{\tilde{B}}(x_i))\right)}{n}\,.


\end{multline*}


  \end{enumerate}


  


  В приведенных формулах $\tilde{A}(x_i)$ и~$\tilde{B}(x_i)$~--- нечеткие 


множества, определенные на базовом множестве признаков $x_i\hm\in X$, 


$i\hm=\overline{1,n}$, которые, в~свою очередь, могут принимать как 


количественные (числовые), так и~качественные (лингвистические) 


значения, обусловливающие применение той или иной формулы. В~первом 


случае целесообразно воспользоваться мерой сходства по Лукасевичу, 


имеющей смысл расстояния между элементами, а во втором~--- мерой 


сходства по Заде, используемой при выполнении операций над логическими 


переменными и~высказываниями.


  


  В случае когда в~описании ситуаций одновременно присутствуют 


и~количественные, и~качественные параметры, целесообразно применять 


комбинированную меру сходства





\noindent


  $$


  C_0\left( \tilde{A},\tilde{B}\right) =\beta m_Z\left(\tilde{A},\tilde{B}\right) 


+(1-\beta) m_L\left(\tilde{A},\tilde{B}\right),


  $$


  в~которой конкретные значения коэффициента $\beta\hm\in [0,1]$ 


выбираются экспертами исходя из содержательного смысла используемых 


параметров~[10].


  


\section{Представление семейной структуры как нечеткого


множества}





  Имеются два множества: множество персон, относящихся к~определенной 


семейной структуре, информация о~которых уже зафиксирована 


в~семантической сети технологии, и~множество персон, упоминаемых 


в~некотором номинативном историческом документе. На их основе 


формируются два нечетких множества с~соответствием, определяемым 


связанностью элементов, определенной в~результате выполнения процедур 


ВИС (P-связностью), и~структурой отношений родства.


  


  Пусть $A$~--- множество узлов семантической сети, соответствующих 


составу семьи в~предыдущий период в~соответствии с~ранее обработанными 


документами семейной структуры. 


  


  


  Пусть $B$~--- множество персон, упоминаемых в~обрабатываемом 


документе семейной структуры, для которых уже в~результате анализа 


документа выявлены отношения родства и~для некоторых из них 


установлены P-свя\-зи с~соответствующими элементами множества~$A$ 


средствами ВИС. Подмножество~$B^0$ множества~$B$, для элементов 


которого установлены P-свя\-зи с~элементами~$A$, будем называть ядром 


множества~$B$. Ему соответствует ядро~$A^0$ множества~$A$. 


Элемент~$B^0$, связанный \mbox{P-связью} с~элементом $a\hm\in A^0$, 


обозначим~$p(a)$. Вероятность соответствия, определенная 


методом ВИС для $a_i\hm\in A^0$, обозначим как~$v_i$.


  


  Как показано выше, дихотомическая оценка несовпадений структур 


слишком груба и~субъективна. Целесообразно использовать экспертную 


оценку по порядковой шкале. Существует ряд методов построения по 


экспертным оценкам функции принадлежности нечеткого множества. 


В~целях снижения субъективного влияния эксперта на результаты 


построения функций принадлежности, которое в~ряде случаев приводит 


к~сдвиганию оценки объектов в~направлении концов оценочной шкалы, 


в~рамках косвенных методов построения функций принадлежности 


используются: метод статистических данных; метод парных сравнений; 


метод экспертных оценок; параметрический метод; метод интервальных 


оценок; метод ранговых оценок~[11]. Целесообразно строить алгоритм 


построения функции принадлежности на базе модифицированного метода 


Ягера и~модифицированного метода анализа иерархий Саати~[12]. Значение 


построенной функции принадлежности для узла~$x_i$ одного графа, не 


имеющего аналога в~другом графе, т.\,е.\ не принадлежащего ядру, 


отражающее оценку степени объяснимости несоответствия для данного узла, 


обозначим как~$o_i$.


  


  Обозначим $\overline{A}\hm= A\cup B\backslash B^0$, $\overline{B}\hm= 


B\cup A\backslash A^0$. Между элементами множеств~$\overline{A}$ 


и~$\overline{B}$ существует взаимно однозначное соответствие (биекция) $f: 


\overline{A}\hm\leftrightarrow \overline{B}$:


  


  \noindent


  $$


  f\left(a_i\in \overline{A}\right)= \begin{cases}


  p(a_i), & a_i\in A^0\,;\\


  a_i, & a_i\in A\backslash A^0\cup B\backslash B^0\,.


  \end{cases}


  $$


    Это позволяет рассмотреть два нечетких множества~$\tilde{A}$ 


и~$\tilde{B}$, определенных на универсальном множестве~$\overline{B}$ 


и~с~функциями принадлежности~$\mu_{\tilde{A}}$ и~$\mu_{\tilde{B}}$ 


соответственно:


  $$


  \mu_{\tilde{A}} (x_i)= \begin{cases}


  1 & \vert f^{-1}(x_i)\in A^0\,;\\


  1& \vert x_i\in A\backslash A^0\,;\\


  o_i & \vert x_i\in B\backslash B^0\,,


  \end{cases}\quad


  \mu_{\tilde{B}} (x_i)= \begin{cases}


  v_i & \vert x_i\in B^0\,;\\


  1& \vert x_i\in B\backslash B^0\,;\\


  o_i & \vert x_i\in A\backslash A^0\,.


  \end{cases}


  $$


  


  Функция $\mu_{\tilde{A}}$ отражает тот факт, что множество~$A$ 


является эталонным, а~элементы~$B$, отсутствующие в~$A$, относятся 


к~нему с~вероятностями, определенными на базе экспертных оценок степени 


объяснимости несоответствия. Функция~$\mu_{\tilde{B}}$ отражает тот 


факт, что элементы ядра множества~$B$ соотносятся с~элементами 


множества~$A$ с~вероятностями, определенными методом ВИС, прочие 


элементы множества~$B$ относятся к~нему безусловно, а~элементы~$A$, 


отсутствующие в~$B$, относятся к~нему с~вероятностями, определенными 


на базе экспертных оценок степени объяснимости несоответствия.





\section{Решение о связывании записей}





  Решение о связывании или несвязывании записей номинативных 


источников семейной структуры целесообразно принимать, оценивая 


значение меры сходства двух нечетких множеств~$\tilde{A}$ и~$\tilde{B}$, 


построенных описанным выше способом. Поскольку функции 


принадлежности строились на основании как экспертных данных, так 


и~результатов расчетов, целесообразно применять комбинированную меру 


сходства


  $$


  C_0\left( \tilde{A},\tilde{B}\right) =\beta m_Z\left(\tilde{A},\tilde{B}\right) 


+(1-\beta) m_L \left( \tilde{A},\tilde{B}\right),


  $$


где $m_Z$~--- мера сходства по Заде;  $m_L$~--- мера сходства по 


Лукасевичу.


  


  Коэффициент~$\beta$, отражающий соотношение количественных 


и~качественных параметров, может быть вычислен по формуле:


  $$


  \beta= \fr{\vert A\backslash A^0 \vert +\vert B\backslash B^0 \vert} {\vert 


A\vert+ \vert B\vert +\vert B\backslash B^0\vert}\,.


  $$





\section{Выводы}





  Предложенный алгоритм позволяет значительно улучшить входящий 


в~состав технологии ПКИИ механизм автоматизации проведения такой 


важной составляющей кон\-крет\-но-ис\-то\-ри\-че\-ско\-го исследования, как 


идентификация персон в~номинативных исторических документах семейной 


структуры, за счет снижения влияния субъективного мнения исследователя 


на результат и~увеличения точности результата за счет отказа от дихотомии 


при классификации несовпадений структур отношений родства.


  


  Дальнейшее развитие предложенного алгоритма возможно за счет 


использования экспертных систем при формировании оценок несовпадений 


этих структур. 


  


{\small\frenchspacing


{%\baselineskip=10.8pt


\addcontentsline{toc}{section}{Литература}


\begin{thebibliography}{99}


  \bibitem{1-ad}


  \Au{ Помникова А.\,Ю.} Семейная история в~дискурсивном пространстве~/\!\!/ Вестник 


Мининского университета, 2019. Т.~7.  №\,1. Ст.~9. 22~с.


  \bibitem{2-ad}


  \Au{Адамович И.\,М., Волков О.\,И.} Технология распределенного автоматизированного 


анализа исторических текстов /\!\!/ Системы и~средства информатики, 2016. Т.~26. №\,3. 


С.~148--161.


  \bibitem{3-ad}


  \Au{Адамович И.\,М., Волков~О.\,И.} Единая технология поддержки  


кон\-крет\-но-ис\-то\-ри\-че\-ских исследований /\!\!/ Системы и~средства информатики, 


2019. Т.~29. №\,1. С.~194--205.


  \bibitem{4-ad}


  \Au{Антонов Д.\,Н.} Восстановление истории семей: метод, источники, анализ:\linebreak Дис.\ 


\ldots\ канд. ист. наук.~--- М.: РГГУ, 2000. 290~с.


  \bibitem{5-ad}


  \Au{Торвальдсен Г.} Номинативные источники в~контексте всемирной истории 


переписей: Россия и~Запад~/\!\!/ Известия Уральского федерального университета. Сер.~2: 


Гуманитарные науки, 2016. Т.~18. №\,3. С.~9--28.


  \bibitem{6-ad}


  \Au{Юрченко Н. \,Л.} Некоторые проблемы использования ревизских сказок как 


источника по исторической демографии~/\!\!/ Вспомогательные исторические 


дисциплины, 1993. T.~XXIV. С.~183--189.


  \bibitem{7-ad}


  \Au{Адамович И.\,М., Волков О.\,И.} Структурный подход к~связыванию записей 


в~технологии поддержки кон\-крет\-но-ис\-то\-ри\-че\-ских исследований~/\!\!/ Системы 


и~средства информатики, 2022. Т.~32. №\,1. С.~94--103.


  \bibitem{8-ad}


  \Au{Орлов А.\,И.} Экспертные оценки.~--- М., 2002. 31~с.


  \bibitem{9-ad}


  \Au{Конышева Л.\,К., Серова~Т.\,А.} Элементы теории нечетких множеств.~--- 


  Екатеринбург: Рос. гос. проф.-пед. ун-т, 2007. 129~с.


  \bibitem{10-ad}


  \Au{Целых А.\,А.} Метод оценки разрыва в~компетенциях выпускника вуза~/\!\!/ 


Фундаментальные исследования, 2013. №\,11-8. С.~1602--1606. 


  \bibitem{11-ad}


  \Au{Каид В.\,А.\,А.} Методы построения функций принадлежности нечетких 


множеств~/\!\!/ Известия ЮФУ. Технические науки, 2013. №\,2(139). С.~144--153. 


  \bibitem{12-ad}


  \Au{Гриценко С.\,А., Храмов В.\,Ю.} Косвенные методы построения функций 


принадлежности сис\-тем поддержки принятия решений с~нечеткой логикой~/\!\!/ Вестник 


ВГУ. Сер.: Сис\-тем\-ный анализ и~информационные технологии, 2017. №\,4. С.~114--125.


       \end{thebibliography}


} }











\vspace*{-3pt}





\hfill{\small\textit{Поступила в~редакцию 28.02.22}}





%\vspace*{8pt}





%\hrule





%\vspace*{2pt}





\newpage





\vspace*{-28pt}





%\hrule





%\vspace*{8pt}





\def\tit{APPROACH TO RECORD LINKING IN~TECHNOLOGY 


OF~CONCRETE HISTORICAL INVESTIGATION SUPPORT BASED~ON~FUZZY~SETS}











\def\titkol{Approach to record linking in~technology 


of~concrete historical investigation support}


% based  on~fuzzy sets}





\def\aut{I.\,M.~Adamovich and O.\,I.~Volkov}





\def\autkol{I.\,M.~Adamovich and O.\,I.~Volkov}





\titel{\tit}{\aut}{\autkol}{\titkol}





\vspace*{-8pt}





\noindent


Federal Research Center ``Computer Science and Control'' of the Russian Academy of Sciences, 44-2~Vavilov Str., 


Moscow 119133, Russian Federation





\def\leftfootline{\small{\textbf{\thepage}


\hfill Sistemy i Sredstva Informatiki~---


Systems and Means of Informatics 2022 vol~32 no\,2}


}%


 \def\rightfootline{\small{Sistemy i Sredstva Informatiki~---


 Systems and Means of Informatics 2022 vol~32 no\,2


\hfill \textbf{\thepage}}}





\vspace*{3pt}


  


  





    


  \Abste{The article is devoted to the further development of the distributed 


technology of concrete historical investigation support based on the principles of 


crowdsourcing and focused on a~wide range of users which are nonprofessional 


historians and biographers. The development is carried out by modifying the 


algorithm for automatic record linking when processing nominative documents of 


family structure based on the Family reconstitution method, successfully used in 


historical demography, and the structure of family relationship graphs. The 


modification consists in relying on the theory of fuzzy sets and using the similarity 


measures to make a~decision about records linking in the case when the algorithms 


of the family reconstitution method do not allow an unambiguous conclusion about 


the linkage. The article describes the methodology for constructing membership 


functions and substantiates the use of a~similarity measure based on the 


Lukasiewicz and Zade similarity measures.}


  


  \KWE{concrete historical investigation; distributed technology; record 


linking; sources of family structure; fuzzy sets}


  


\DOI{10.14357\!/\!08696527220213} 





%\vspace*{-16pt}








 %\Ack


%\noindent





 








%\vspace*{-6pt}





\renewcommand{\bibname}{\protect\rmfamily References}


%\renewcommand{\bibname}{\large\protect\rm References}





{\small\frenchspacing


{%\baselineskip=10.8pt


\addcontentsline{toc}{section}{References}


\begin{thebibliography}{99}


  \bibitem{1-ad-1}


\Aue{Pomnikova, A.\,Yu.} 2019. Semeynaya istoriya v~diskursivnom prostranstve 


[Family stories in different types of discourse]. \textit{Vestnik of Minin University} 


7(1):9. 22~p.


  \bibitem{2-ad-1}


  \Aue{Adamovich, I.\,M., and O.\,I.~Volkov.} 2016. Tekhnologiya 


raspredelennogo avtomatizirovannogo analiza istoricheskikh tekstov [The 


distributed automated technology of historical texts analysis]. \textit{Sistemy 


i~Sredstva Informatiki~--- Systems and Means of Informatics} 26(3):148--161.


\bibitem{3-ad-1}


\Aue{Adamovich, I.\,M., and O.\,I.~Volkov.} 2019. Edinaya tekhnologiya 


podderzhki konkretno-istoricheskikh issledovaniy [Unified technology of concrete 


historical research support]. \textit{Sistemy i~Sredstva Informatiki~--- Systems and 


Means of Informatics} 29(1):194--205.


  \bibitem{4-ad-1}


\Aue{Antonov, D.\,N.} 2000. Vosstanovlenie istorii semey: metod, istochniki, 


analiz [Family reconstitution: Method, sources, and analysis].  Moscow: 


RGGU. PhD Thesis. 290~p.


  \bibitem{5-ad-1}


\Aue{Thorvaldsen, G.} 2016. Nominativnye istochniki v~kontekste vsemirnoy istorii 


perepisey: Rossiya i~Zapad [Nominative data and global census


history: Russia and the West]. \textit{Izvestiya Ural'skogo 


federal'nogo universiteta. Ser.~2: Gumanitarnye nauki} [Izvestia. Ural Federal 


University~Journal. Ser.~2. Humanities and Arts] 18(3):9--28.


  \bibitem{6-ad-1}


\Aue{Yurchenko, N.\,L.} 1993. Nekotorye problemy ispol'zovaniya revizskikh 


skazok kak istochnika po istoricheskoy demografii [Some problems of using 


census lists as a~source for historical demography]. \textit{Vspomogatel'nye 


istoricheskie distsipliny} [Auxiliary Historical Disciplines] XXIV:183--189.


  \bibitem{7-ad-1}


\Aue{Adamovich, I.\,M., and O.\,I.~Volkov.} 2022. Strukturnyy podkhod 


k~svyazyvaniyu zapisey v~tekhnologii podderzhki konkretno-istoricheskikh 


issledovaniy [Structural approach to record linking in technology of concrete 


historical investigation support]. \textit{Sistemy i~Sredstva Informatiki~--- Systems 


and Means of Informatics} 32(1):94--103. 


  \bibitem{8-ad-1}


  \Aue{Orlov, A.\,I.} 2002.


\textit{Ekspertnye otsenki} [Expert assessments]. Moscow. 31~p. Available at: {\sf 


http://www.aup.ru/books/m154/} (accessed March~10, 2022).


  \bibitem{9-ad-1}


\Aue{Konysheva, L.\,K., and T.\,A.~Serova.} 2007. \textit{Elementy teorii 


nechetkikh mnozhestv} [Fuzzy-set theory elements]. Ekaterinburg: RGPPU. 129~p.


  \bibitem{10-ad-1}


\Aue{Tselykh, A.\,A.} 2013. Metod otsenki razryva v~kompetentsiyakh 


vypusknika vuza [Method for competency gap analysis of a~university graduate]. 


\textit{Fundamental'nye issledovaniya} [Fundamental Research] 11-8:1602--1606.


  \bibitem{11-ad-1}


\Aue{Qaid, W.\,A.\,A.} 2013. Metody postroeniya funktsiy prinadlezhnosti 


nechetkikh mno\-zhestv [Methods construction membership function of fuzzy sets]. 


\textit{Izvestiya YuFU. Tekhnicheskie nauki} [Izvestiya SFedU: Engineering Sciences] 


2(139):144--153.


  \bibitem{12-ad-1}


  \Aue{Gritsenko, S.\,A., and V.\,U.~Khramov.} 2017. Kosvennye metody 


postroeniya funktsiy prinadlezhnosti sistem podderzhki prinyatiya resheniy 


s~nechetkoy logikoy [Indirect methods of constructing the functions of accessories 


of decision support systems with fuzzy logic]. \textit{Vestnik VGU. Ser.: Sistemnyy 


analiz i~informatsionnye tekhnologii} [Proceedincs of Voronezh State University. 


Ser.: Systems Analysis and Information Technologies] 4:114--125.


  \end{thebibliography}


} }





\vspace*{-3pt}





\hfill{\small\textit{Received February 28, 2022}} 





%\vspace*{-14pt} 


  


  \Contr


  


  \noindent


  \textbf{Adamovich Igor M.} (b.\ 1934)~--- Candidate of Science (PhD) in 


technology, leading scientist, Institute of Informatics Problems, Federal Research 


Center ``Computer Science and Control'' of the Russian Academy of Sciences,  


44-2~Vavilov Str., Moscow 119133, Russian Federation; 


\mbox{Adam@amsd.com}


  


  \vspace*{3pt}


  


  \noindent


  \textbf{Volkov Oleg I.} (b.\ 1964)~--- leading programmer, Institute of 


Informatics Problems, Federal Research Center ``Computer Science and Control'' 


of the Russian Academy of Sciences, 44-2~Vavilov Str., Moscow 119133, Russian 


Federation; \mbox{Volkov@amsd.com}





  


\label{end\stat}





\renewcommand{\bibname}{\protect\rm Литература} 


